\begin{abstract}
A procedure can be correct where it was learned yet incomplete where it is reused. We study coding agents supplied with procedural memory in target tasks that change a security relevant source assumption. Thirteen synthetic families cross four memory conditions and two repetitions in six model and agent configurations, yielding \RecordedRuns{} recorded runs and \ValidRuns{} technically valid outcomes. The primary endpoint is functional completion that fails a focal security witness. Estimated changes from supplying correct memory range from $\PrimaryURangeLow$ to $\PrimaryURangeHigh$ percentage points across configurations; every corresponding bootstrap interval includes zero. Lower endpoint values sometimes accompany lower functionality. A generic applicability reminder lowers the endpoint in all three Codex configurations, including under bounds for missing outcomes in the fixed recorded cohort; the MiniSWE configurations do not share that pattern. In a 52-run sample selected by a fixed rule, two model assisted coding passes agree on \CodeTBothYes{} implementations with a target protection but only \CodeABothYes{} with explicit assumption recognition in the visible trace before editing. Matched transcripts and patches show unchecked reuse, successful adaptation, and distinct mechanisms behind the same measured failure. These observations distinguish three outcomes of procedural transfer that a single success metric can conflate: adaptation, functional completion with a focal violation, and failure to complete.
\end{abstract}
